% ============================================================
% EJERCICIO 4 — Corrección asistida por IA (Módulo 2)
% Este documento COMPILA pero tiene al menos 8 problemas de buenas prácticas.
% Tareas: 1) Pida a una IA que los detecte (pegue este código y el prompt).
%         2) NO acepte ciegamente: compruebe cada sugerencia compilando.
%         3) Entregue ejercicio4_corregido.tex y mejoras.txt con las
%            sugerencias de la IA que implementó y las que RECHAZÓ (y por qué).
% ============================================================
\documentclass[12pt]{article}
\usepackage[utf8]{inputenc}
\usepackage{amsmath}
\usepackage{graphicx}

\title{Ley de Ohm}
\author{Estudiante}
\begin{document}
\maketitle
\tableofcontents

\section{Introduccion}
La ley de Ohm dice que V = I x R.\\
\\
Para dos resistencias en serie: $$R_{eq} = R_1 + R_2$$

\section{Resultados}
Los resultados se ven en la Tabla 1 y en la Figura 1.

\begin{table}[h]
\begin{tabular}{|c|c|}
\hline
V (V) & I (mA) \\
\hline
1.0 & 9.8 \\
5.0 & 49.2 \\
\hline
\end{tabular}
\end{table}

\fbox{Figura aquí}

\end{document}
